\subsection{Antecedentes Investigativos}

La selección de estudios incluidos en la presente revisión se realizó aplicando de
forma sistemática criterios de inclusión y exclusión definidos a partir de los
objetivos específicos del proyecto y de las variables de investigación. Cada estudio
fue evaluado en dos etapas: revisión de título y resumen, seguida de lectura completa
del texto en los casos que superaron el primer filtro. Los criterios aplicados se
detallan en la Tabla~\ref{tab:criterios-ant}.

\begin{table}[H]
\centering
\small
\renewcommand{\arraystretch}{1.3}
\begin{tabular}{|p{0.22\textwidth}|p{0.33\textwidth}|p{0.33\textwidth}|}
\hline
\textbf{Criterio} & \textbf{Inclusión} & \textbf{Exclusión} \\
\hline
Vigencia &
  2016--2026 (2019--2026 para chatbots/agentes conversacionales) &
  Anterior a ese rango, salvo estudio seminal justificado \\
\hline
Accesibilidad &
  Texto completo disponible &
  Fuente restringida o solo resumen \\
\hline
Tipo de fuente &
  Revista arbitrada, congreso arbitrado, tesis en repositorio institucional &
  Blogs, notas de prensa, contenido comercial o ejemplos ficticios \\
\hline
Arbitraje &
  Revisado por pares &
  Preprint sin revisión por pares, salvo excepción justificada \\
\hline
Idioma &
  Español o inglés &
  Otro idioma sin traducción disponible \\
\hline
Duplicidad &
  Se conserva una sola versión del mismo estudio &
  Versión duplicada de un estudio ya incluido \\
\hline
Redundancia &
  Se conserva el de mayor aporte entre estudios equivalentes &
  Repite hallazgo o enfoque ya cubierto sin aportar nada nuevo \\
\hline
Relevancia temática mínima &
  Toca flujo clínico-administrativo, gestión y eficiencia en salud,
  integración de sistemas, agente conversacional o medición de tiempos &
  No toca ninguno de esos ejes temáticos \\
\hline
Verificabilidad &
  Referencia (DOI, autores) coincide con el contenido &
  Atribución dudosa o inconsistente \\
\hline
\end{tabular}
\caption{Criterios de inclusión y exclusión aplicados en la selección de antecedentes}
\label{tab:criterios-ant}
\end{table}

La aplicación de los criterios de inclusión y exclusión permitió filtrar un conjunto
inicial de aproximadamente 122 artículos científicos identificados en las bases de
datos consultadas. Sobre ese conjunto depurado se aplicó, en una segunda etapa, la
matriz de pertinencia descrita en la Tabla~\ref{tab:pertinencia}, la cual evaluó cada
estudio en cinco criterios vinculados directamente a los objetivos específicos y las
variables de investigación, con una escala de tres niveles:
\textbf{--1}~No pertinente, \textbf{0}~Medianamente pertinente,
\textbf{+1}~Altamente pertinente.

\begin{table}[H]
\centering
\small
\renewcommand{\arraystretch}{1.3}
\begin{tabular}{|p{0.22\textwidth}|p{0.22\textwidth}|p{0.22\textwidth}|p{0.22\textwidth}|}
\hline
\textbf{Criterio} & \textbf{--1 No pertinente} & \textbf{0 Med. pertinente} & \textbf{+1 Alt. pertinente} \\
\hline
Alineación con OE &
  No aborda flujo, agente conversacional ni implementación/medición de tiempos &
  Toca el dominio general (salud digital/TI en salud) sin responder a un OE específico &
  Sustenta directamente OE1 (flujo/requerimientos), OE2 (arquitectura del agente)
  u OE3 (implementación/medición) \\
\hline
Alineación con VI/VD &
  Sin relación con estandarización/integración ni con tiempo de gestión &
  Las menciona tangencialmente, sin medirlas &
  Sustenta explícitamente la VI y/o la VD con evidencia o medición concreta \\
\hline
Contexto de aplicación &
  Ajeno a salud/servicios con agendamiento &
  Salud en general, no odontológico/administrativo específico &
  Odontológico, o gestión de citas/administrativa directamente comparable \\
\hline
Aporte técnico/empírico &
  Descriptivo o divulgativo, sin datos ni diseño aplicable &
  Marco conceptual o taxonomía útil como referencia teórica &
  Arquitectura, datos cuantitativos o método directamente reutilizable \\
\hline
Utilidad para problema/ justificación &
  No aporta evidencia de magnitud ni de pertinencia de la solución &
  Evidencia parcial o indirecta &
  Evidencia sólida de magnitud del problema o beneficio de la solución \\
\hline
\end{tabular}
\caption{Escala de pertinencia aplicada en la selección de antecedentes
  (\textbf{+1}~Altamente pertinente; \textbf{0}~Medianamente pertinente;
   \textbf{--1}~No pertinente)}
\label{tab:pertinencia}
\end{table}

Tras aplicar ambas etapas de filtrado, quedaron seleccionados los trece estudios con
mayor puntuación de pertinencia. La distribución de esas referencias según la fuente
de publicación se presenta en la Tabla~\ref{tab:fuentes-ant}.

\begin{table}[H]
\centering
\renewcommand{\arraystretch}{1.3}
\begin{tabular}{|l|c|}
\hline
\textbf{Base de datos / Fuente} & \textbf{Artículos seleccionados} \\
\hline
JMIR Publications                            & 3 \\
\hline
Science Direct                               & 2 \\
\hline
Springer (Journal of Biomedical Informatics) & 1 \\
\hline
IEEE Xplore / IJATCSE                        & 1 \\
\hline
International Journal of Computing           & 1 \\
\hline
Suranaree Journal of Social Science          & 1 \\
\hline
IJCRT                                        & 1 \\
\hline
Revista Cubana de Informática Médica         & 1 \\
\hline
Dominio de las Ciencias (Dom.\ Cien.)        & 1 \\
\hline
Iustitia Socialis (Dialnet)                  & 1 \\
\hline
Repositorio UCE                              & 1 \\
\hline
Procedia Computer Science                    & 1 \\
\hline
\textbf{Total}                               & \textbf{13} \\
\hline
\end{tabular}
\caption{Distribución de artículos seleccionados por fuente de publicación}
\label{tab:fuentes-ant}
\end{table}

Con base en este conjunto de referencias recopiladas, se realiza el análisis
cualitativo de la literatura, organizando los hallazgos de acuerdo con los objetivos
específicos del proyecto. Para desarrollar el primer objetivo (OE1), centrado en el
levantamiento de requerimientos funcionales, el análisis de normativas y el diagnóstico
del flujo clínico y administrativo actual, es necesario revisar investigaciones previas
sobre las metodologías y herramientas utilizadas para modelar procesos en el sector
salud, así como los factores normativos y organizacionales que determinan la calidad del
registro clínico en contextos de recursos limitados. Para el segundo objetivo (OE2),
referido al diseño de la arquitectura del agente conversacional y del sistema de
información, se examinan antecedentes de arquitecturas de software basadas en
microservicios, estándares de interoperabilidad y patrones de integración de chatbots
en plataformas de mensajería. Para el tercer objetivo (OE3), orientado a la
implementación y evaluación de impacto, se revisan estudios que aplican metodologías
de tiempos y movimientos, diseños antes-después y esquemas de pruebas funcionales en
entornos clínicos.

% -------------------------------------------------------
% PÁRRAFOS POR ESTUDIO
% -------------------------------------------------------

En \cite{gomes2018bpm} se desarrolló un caso de estudio cualitativo en un hospital
de gran escala en Portugal, donde se aplicó la notación BPMN sobre la plataforma de
historia clínica electrónica AIDA para modelar cuatro módulos del sistema: ambulatorio,
consultas, hospitalización y urgencias. Se identificaron tres niveles de problemas de
gestión de procesos —estratégico, táctico y operativo— entre los que destacan la falta
de estándares, debilidades en la especificación de los procesos y brechas entre el
diseño y la ejecución real. El estudio no llega a la fase de implementación ni reporta
cifras de impacto, por lo que sus aportaciones son de carácter conceptual; sin embargo,
evidencia que la ausencia de formalización de procesos clínico-administrativos es una
barrera reconocida en la literatura, independientemente de la escala institucional.

En \cite{staras2021using} se presentó un tutorial metodológico estructurado en cuatro
pasos para analizar el flujo de trabajo clínico antes de implementar una intervención
de eHealth, validado en cuatro clínicas pediátricas y de medicina familiar en Estados
Unidos. Los pasos propuestos son: (i)~identificación de componentes discretos del
flujo, (ii)~evaluación mediante observación directa con lista de verificación
estandarizada y marcación de tiempo, (iii)~triangulación con entrevistas y diagramas
de flujo, y (iv)~propuesta de estrategias de implementación con los propios usuarios
clínicos. El caso de estudio involucró trece observaciones de visitas de pacientes,
registrando tiempos de espera de hasta 59 minutos en sala (media 12{,}6 min,
DE 16{,}1~min) y hasta 25 minutos sin contacto con personal antes de la consulta.
Se halló que realizar el diagnóstico del flujo antes de la implementación permite
anticipar barreras operativas e incrementa la aceptación del personal.

En \cite{tudorcar2020conversational} se llevó a cabo una revisión sistemática de
alcance (\textit{scoping review}) que mapeó las aplicaciones, vacíos de investigación
y desafíos de los agentes conversacionales en la atención médica a partir de
47~reportes empíricos. Se identificó que la mayoría de los chatbots revisados emplean
procesamiento de lenguaje natural, operan sobre texto y se implementan a través de
aplicaciones móviles. Un hallazgo técnico destacado es que los chatbots integrados en
plataformas de mensajería masiva —frente a aplicaciones propietarias— presentan menores
tasas de abandono. Se señaló como vacío crítico que ninguno de los estudios revisados
midió si la tecnología conversacional mejora la productividad del personal o reduce la
carga de trabajo en comparación con el desempeño estándar del personal administrativo.

En \cite{wongperez2023hl7} se realizó una revisión sistemática de la literatura
aplicando la metodología de Petersen en cinco pasos, con el objetivo de determinar el
estado del arte y las tendencias de utilización del estándar HL7 en sistemas de
información sanitarios. La búsqueda en IEEE Xplore y ScienceDirect recuperó
inicialmente 428\,569 entradas; tras los filtros de selección, se incluyeron
21~artículos del período 2017--2021, con el 42{,}8\,\% provenientes de Estados Unidos.
Se documentó que HL7~FHIR fue la variante de mayor aplicación en los años analizados,
y se destacó como hallazgo particular un antecedente arquitectónico que combina
microservicios, HL7~FHIR y AIML para un chatbot de pacientes crónicos, integración
que se presentó como técnicamente sólida y replicable.

En \cite{roca2020microservice} se propuso y validó mediante un prototipo
(\textit{proof of concept}) una arquitectura de chatbot basada en microservicios para
el soporte a pacientes crónicos, sustentada en tres pilares: escalabilidad mediante
microservicios independientes desplegados en contenedores Docker, modelo estándar de
datos mediante HL7~FHIR, y modelado estándar de conversación mediante AIML~2.0. La
arquitectura se organiza en un microservicio Proxy, un API Gateway y microservicios
funcionales (cuestionarios, citas, registro, especialista), comunicados vía JSON/HTTPS.
Como caso de extensión se incorporaron tres nuevos microservicios sin modificar el
núcleo del sistema, demostrando la modularidad del diseño. También se desarrolló un
generador automático de archivos AIML a partir de recursos FHIR tipo
\textit{Questionnaire}.

En \cite{alejandrino2023dental} se presentó el plan de desarrollo de un Sistema de
Información para clínicas dentales privadas con integración de un chatbot basado en
reglas y procesamiento de lenguaje natural (NLP), en el contexto de una clínica dental
privada en Filipinas. El sistema propuesto contempla módulos de administración,
facturación, inventario, pacientes y reportes, más un módulo de chatbot para
agendamiento de citas, desarrollado sobre el \textit{framework} Laravel con
arquitectura web de tres capas. A partir de análisis de procesos y entrevistas con
\textit{stakeholders}, se proyectaron reducciones en tiempos administrativos para
cinco procesos clave, con una disminución estimada del 68{,}75\,\% en el tiempo total
(de 80 a 25~minutos). Estas cifras corresponden a proyecciones de un plan de proyecto
—no a mediciones empíricas post-implementación— y deben interpretarse con cautela.

En \cite{arcila2022microservice} se diseñó y validó mediante prototipo una arquitectura
de software basada en microservicios para mejorar la disponibilidad de registros de
salud dental en Perú. La arquitectura se compone de cuatro microservicios —Paciente,
Registro Médico Dental, Odontograma y Proveedor de Servicios Dentales— comunicados
mediante HTTP REST a través de un API Gateway con autenticación por JSON Web Token
(JWT), desplegados en Microsoft Azure con contenedores independientes. Las pruebas de
carga con JMeter demostraron que con una instancia el sistema soporta 21~solicitudes
por segundo con 100\,\% de disponibilidad, y con tres instancias alcanza
63~solicitudes por segundo en las mismas condiciones, verificando los atributos de
calidad definidos bajo el estándar ISO/IEC~25010.

En \cite{hmadchayyapum2025line} se diseñó y desarrolló un chatbot sobre la plataforma
LINE para la reserva y consulta de servicios odontológicos de un centro de salud oral
universitario en Tailandia, aplicando el ciclo de vida de desarrollo de sistemas (SDLC)
en cuatro etapas. El levantamiento de requisitos se realizó mediante entrevistas en
profundidad a tres funcionarios (muestreo intencional) y cinco usuarios previos
(muestreo accidental). Se implementaron cinco funciones: reservar turno, consultar
cola, visualizar turno propio, cancelar y realizar consultas. La evaluación de
usabilidad con quince usuarios obtuvo una media de 4{,}20 sobre 5{,}00 (DE~0{,}37),
con la facilidad de uso como dimensión mejor valorada. No se midió la reducción del
tiempo administrativo del personal.

En \cite{hindelang2024chatbots} se sintetizó la evidencia de 18~estudios
—15~observacionales y 3~ensayos clínicos— sobre chatbots aplicados a la toma de
historia clínica, en una revisión sistemática multi-institucional. Las tecnologías
identificadas incluyen NLP, \textit{machine learning}, reconocimiento de voz, árboles
de decisión, sistemas basados en reglas, AIML y plataformas como IBM Watson, así como
canales web, tabletas y Telegram. Se reportó un ahorro de tiempo del 57{,}3\,\% en un
estudio de anamnesis asistida por chatbot, completitud del 73{,}3\,\% en historias
clínicas recopiladas por agente conversacional, y una puntuación de usabilidad de~96.
Se identificaron como barreras la resistencia de usuarios, la privacidad y el
desempeño insuficiente en situaciones clínicas complejas.

En \cite{shaikh2026smartmedibook} se desarrolló e implementó un sistema web de reserva
de citas hospitalarias con inteligencia artificial en India, mediante una arquitectura
que combina frontend en React.js, backend en Node.js/Express.js, base de datos
MongoDB Atlas desplegada en la nube y un módulo de chatbot independiente desarrollado
con FastAPI, comunicados mediante REST APIs. El sistema implementa autenticación con
JWT, control dinámico de disponibilidad de franjas horarias y prevención de doble
reserva. La validación funcional reportó cinco casos de prueba (TC01--TC05) con
resultado \textit{Pass} en su totalidad. No se reportaron cifras cuantitativas
pre-post del tiempo administrativo por paciente.

En \cite{laraespinoza2026gestion} se realizó un estudio de tipo mixto —descriptivo,
explicativo y propositivo— con diseño no experimental transversal, mediante
45~encuestas y 6~entrevistas al personal de un centro de salud en Ambato, Ecuador,
con el objetivo de proponer un manual de procedimientos para la automatización de la
historia clínica electrónica. Se encontró que el 66{,}7\,\% del personal aún utiliza
registros en papel, el 75{,}6\,\% señala la ausencia de un sistema híbrido funcional,
el 80\,\% considera la infraestructura tecnológica inadecuada y el 75{,}6\,\% no ha
recibido capacitación digital en los últimos doce meses. El manual propuesto fue
validado mediante juicio de expertos con escala Likert, obteniendo una viabilidad del
94\,\% (141 de 150 puntos posibles).

En \cite{martinez2022proteccion} se analizó, mediante un enfoque cualitativo-jurídico,
el nivel de cumplimiento de la Ley Orgánica de Protección de Datos Personales (LOPDP)
del Ecuador respecto a los datos de salud, tomando como caso la historia clínica
electrónica de un subsistema del Sistema Nacional de Salud ecuatoriano. Se examinó el
Informe de Contraloría DNAI-AI-0050-2017, que documentó la ausencia de procedimientos
para administración de usuarios, identificación única de accesos, acuerdos de
confidencialidad y pistas de auditoría identificables. Se identificaron brechas entre
la normativa vigente —artículos 10 y 31 de la LOPDP, el Acuerdo Ministerial 1190
sobre HL7 y el Reglamento de Historia Clínica Electrónica (Acuerdo 0009-2017)— y la
implementación técnica real, con incumplimientos en interoperabilidad, calidad de
datos y protocolos de anonimización.

En \cite{uce2016caracteristicas} se evaluaron 293~formularios 033 de historia clínica
única odontológica del Ministerio de Salud Pública del Ecuador, completados por
estudiantes durante el período septiembre~2015--febrero~2016, mediante un estudio
observacional analítico y transversal. Con diecisiete indicadores de calidad de
registro y cinco de características físicas, y análisis estadístico Chi-cuadrado
(SPSS, IC~95\,\%), se encontró que menos del 50\,\% de los indicadores clave
—incluyendo diagnóstico, plan terapéutico, índice CPOD, antecedentes y motivo de
consulta— fue correctamente completado. Las características físicas de los expedientes
resultaron aceptables, aunque con riesgo de deterioro documentado.